% $Id: zimpl.tex,v 1.8 2002/10/29 10:42:47 bzfkocht Exp $
%* * * * * * * * * * * * * * * * * * * * * * * * * * * * * * * * * * * * * * */
%*                                                                           */
%*   File....: zimpl.tex                                                     */
%*   Name....: Zuse Institute Mathematical Programming Language              */
%*   Author..: Thorsten Koch                                                 */
%*   Copyright (C) 2001 by Author, All rights reserved                       */
%*                                                                           */
%* * * * * * * * * * * * * * * * * * * * * * * * * * * * * * * * * * * * * * */
%
\documentclass[12pt]{article}
\usepackage{times}
\usepackage{a4}
\usepackage[latin1]{inputenc}
\usepackage{amsmath}
\usepackage{amssymb}
\usepackage{xspace}
\usepackage{epsfig}
\usepackage{fancyhdr}
\usepackage{xspace}
\usepackage{url}
%\usepackage{zibtitlepage}
\usepackage[notoday,nofancy]{rcsinfo}
\usepackage[dvips,%
bookmarks,pdffitwindow,pdfcenterwindow=true,pdfstartview=Fit]{hyperref}
\hypersetup{%
pdftitle={ZIMPL User Guide},
pdfsubject={Zuse Institute Mathematical Programming Language},
pdfauthor={Thorsten Koch},
pdfkeywords={Mathematical Modelling Language,Mathematical Programming,Optimization, Algebraic Modelling Language}
}
\rcsInfo $Id: zimpl.tex,v 1.8 2002/10/29 10:42:47 bzfkocht Exp $
%
\parindent0ex
\setcounter{tocdepth}{1}
\setcounter{secnumdepth}{1}
\newcommand{\zimpl}{{\sc Zimpl}\xspace}
\newcommand{\lp}{{\sc lp}\xspace}
\newcommand{\mip}{{\sc mip}\xspace}
\newcommand{\mps}{{\tt mps}\xspace}
\newcommand{\lpf}{{\tt lp}\xspace}
\newcommand{\ibm}{{\sc ibm}\xspace}
\newcommand{\zpl}{{\sc zpl}\xspace}
\newcommand{\ampl}{{\sc ampl}\xspace}
\newcommand{\ilog}{{\sc ilog}\xspace}
\newcommand{\cplex}{{\sc cplex}\xspace}
\newcommand{\code}[1]{{\tt #1}\xspace}
\newcommand{\N}{\ensuremath{\mathbb{N}}}
\newcommand{\Z}{\ensuremath{\mathbb{Z}}}
\headheight5mm
\renewcommand{\footrulewidth}{\headrulewidth}
\lhead{\zimpl}
\chead{}
\rhead{}
\cfoot{\thepage}
\pagestyle{fancy}
%
\begin{document}
%\ZTPAuthor{Thorsten Koch}
%\ZTPTitle{\zimpl User Guide}
%\ZTPInfo{Best preprint of the year 2000}
%\ZTPNumber{01-20}
%\ZTPMonth{August}
%\ZTPPreprint
%\ZTPYear{2001}
%\zibtitlepage

\title{
%\vspace*{-3cm}\epsfig{file=ziblogo2.eps,width=3cm}\\[\bigskipamount]
\LARGE\zimpl\\
\normalsize (Zuse Institute Mathematical Programming Language)}
\author{Thorsten Koch}
\date{\small for Version 1.04\\20. October 2002}
\maketitle
\vfill
\begin{abstract}
\zimpl is a little language to translate the mathematical model of a 
problem into a linear or (mixed-)integer mathematical program
expressed in \lpf or \mps file format which can be read
by a \lp or \mip solver.
\end{abstract}
%\tableofcontents
%\newpage
% -----------------------------------------------------------------------------
% --- Introduction
% -----------------------------------------------------------------------------
\begin{flushright}
{\em May the source be with you, Luke!}
\end{flushright}
\section{Introduction}
Most of the things in \zimpl (and a lot more) can be found in 
the wonderful book about the modelling language \ampl 
from Robert Fourer, David N. Gay and Brian W. Kernighan
\cite{FGK93}.
Indeed if not the guys at \ilog had needed more than three months just
to tell me the price of \ampl for \cplex, I would probably use
\ampl today.
On the other hand, having the source code of a program has its atvantages. The
possibility to run it regardless of architecture and operating system, the
ability to modify it to suite the needs and not having to hassle with license
managers may make a much less powerful program the better choice.
And so \zimpl came into being.

\medskip
A linear program (\lp) might look like this:
$$
\begin{array}{rll}
\min& 2 x + 3 y\\
\mbox{subject to}& x + y& \leq 6\\
&x,y&\ge 0\\
\end{array}
$$
The usual format to store the description of such a problem is \mps
invented by \ibm \cite{IBM97} long ago. 
Nearly all available \lp and \mip solvers can read this format.
While \mps is a nice format to punch into a punch card and at least a
reasonable format to read for a computer it is quite unreadable
for humans. 
\newpage
{
\small
\begin{verbatim}
   NAME        ex1.mps
   ROWS
    N  OBJECTIV          
    L  c0                
   COLUMNS
       x         OBJECTIV             2
       x         c0                   1
       y         OBJECTIV             3
       y         c0                   1
   RHS
       RHS       c0                   6
   ENDATA
\end{verbatim}
}
Another possibility is the \lpf format \cite{ILO00}, which is more
readable but is only supported by a few solvers. 
{
\small
\begin{verbatim}
   Minimize
    cost:  +2 x +3 y
   Subject to
    c0:  +1 x +1 y <= 6
   End
\end{verbatim}
}
But since each coefficient of the matrix $A$ must be stated
explicitly it is also not a desirable choice to develop a mathematical
model.

\medskip
Now, with \zimpl it is possible to write this:
{\small
\begin{verbatim}
   var x;
   var y;
   minimize cost: 2 * x + 3 * y;
   subto c0: x + y <= 6;
\end{verbatim}
}
and have it automatically translated into \mps or \lpf format.
While this looks not much different from what is in the \lpf format,
the difference can be seen, if we use indexed variables.
Here is an example. This is the \lp 
$$
\begin{array}{rl}
\min& 2 x_1 + 3 x_2 + 1.5 x_3\\
\mbox{subject to}&\sum^3_{i=1} x_i \leq 6\\
&x_i\ge 0\\
\end{array}
$$
And this is how to tell it \zimpl
{
\small
\begin{verbatim}
set I      := { 1 to 3 };
param c[I] := <1> 2, <2> 3, <3> 1.5;
var   x[I] >= 0;
minimize value: sum <i> in I : c[i] * x[i];
subto     cons: sum <i> in I : x[i] <= 6;
\end{verbatim}
}

% -----------------------------------------------------------------------------
% --- Invocation
% -----------------------------------------------------------------------------
\section{Invocation}

To run \zimpl on the model given in the file \code{ex1.zpl} type the command
\begin{verbatim}
   zimpl ex1.zpl
\end{verbatim}

The general case is \code{zimpl} \code{[options]} $<$\code{input-files}$>$.
It is possible to give more than one input file. They are read one
after the other and then executed as if they were all in one big file.
The result of running \zimpl will be either an error message or two files.

The two files are the problem generated from the model in either 
\lpf or \mps format with extension \code{.lp} or \code{.mps} and a file
with the extension \code{.tbl} which lists all variable and constraint 
names used in the model and there corresponding name in the problem
file. The reason for this name translation is that the \mps format
can only handle names up to eight characters long. And in \lp format
the length of the names is also restricted to 16 characters.

\medskip
The following options are possible (only the first two are normally
of interest):

\medskip
\begin{tabular}{lp{115mm}}
%\hline
-t {\em format}& Selects the output format. Can be either \code{lp}
which is default or \code{mps}.\\
-o {\em name} & Sets the base-name for the output files.\\ 
& Defaults to the first input file striped of the extension.\\ 
-n {\em cform}& Select the format for the generation of constraint
names. Can be either {\em cm} which will number them $1\ldots n$ with
a `c' in front. {\em cn} will use the name supplied in the
\code{subto} statement and number them $1\ldots n$ within the
statement. {\em cf} will use the name of given with the \code{subto},
then a $1\ldots n$ number like in \emph{cm} and then append all the
local variables from the forall statements.\\
%\hline
-b&Enables bison debugging output.\\
-d&Enables zimpl debugging output.\\
-f&Enables flex debugging output.\\
-h&Prints a usage message.\\
-v&Enables verbose mode.\\
%\hline
\end{tabular}

\medskip
A typical invocation is for example:
\begin{verbatim}
   zimpl -o hardone -t mps data.zpl model.zpl
\end{verbatim}
This reads the files \code{data.zpl} and \code{model.zpl} and
produces \code{hardone.mps} and \code{hardone.tbl}.

% -----------------------------------------------------------------------------
% --- Format
% -----------------------------------------------------------------------------
\section{Format}

Each \zpl-file consists of five types of statements: Sets, parameters,
variables, an objective and constraints.
Each statement ends with a semicolon (\code{;}). 
Everything from a number-sign (\code{\#}) to 
the end of the line is treated as a
comment and is ignored.

If a line starts with the word \code{include} followed by a filename in double
quotation marks, this file is read instead of this line.

\medskip
\zimpl works on the lowest level with two types of data: Strings and
numbers. A string is delimited by double quotation marks (\code{"}).
A number can be given as 2, -6.5 or 5.234e-12. 
Wherever a number or string is required, it is also possible to give a
parameter of the right value type. Also numeric expressions are
allowed instead of just a number. 
The precedence of numeric, boolean and set operators should be the
usual one. If in doubt use parenthesis to be save.

\subsubsection{Numeric expressions}
The following functions are currently implemented: 
$^\wedge$, \code{+}, \code{-}, \code{*}, \code{/}, \code{mod},
\code{div}, \code{abs}, \code{floor}, \code{ceil}, 
\code{log} (logarithm with base 10),
\code{ln} (logarithm with base $e$), \code{exp} ($e^x$), \code{$!$}. 
$^\wedge$ stands for \emph{to the
power of}. \code{mod} is the modulo function and
\code{div} gives the integer part of a division. \code{**} can be used
as synonym for $^\wedge$.
With $\min$ and $\max$ it is possible to find the minimum/maximum
member of an one dimensional set of numeric values.
\code{card} gives the cardinality of a set.

The following is also a numeric expression:

\code{if} \emph{boolean-expression} \code{then}
\emph{numeric-expression} \code{else} \emph{numeric-expression} \code{end}.

\subsubsection{Boolean expressions}
For numbers $<,<=,==,!=,>=,>$ are defined. For strings only $==$ and
$!=$ are available. The expression \emph{tuple} \code{in}
\emph{set-expression} can be used to test set
membership of a tuple. Combinations of boolean-expressions with \code{and},
\code{or} and negation with \code{not} are possible.

\subsection{Sets}
The elements of a set are tuples. Each tuple of a set has the same
number of components. 
The components are either numbers or strings. The type of the n-th
component of each tuple must be the same. A tuple starts and ends with
$<$ and $>$, resp. The components are separated by commas.
If tuples are one-dimensional, it is possible 
to omit the braces in a list of elements, but then 
they must be omitted from all tuples of the set.

\subsubsection{Set functions}
\begin{tabular}{llp{5cm}}
\hline
A cross B   &Cross product & $\{(x,y)| x\in A\wedge y\in B\}$\\
A union B   &Union & $\{x | x\in A\vee x\in B\}$\\
A inter B   &Intersection & $\{x |x\in A\wedge x\in B\}$\\
A without B &Difference & $\{ x |x\in A\wedge x\not\in B\}$\\
A symdiff B &Symmetric difference & 
$\{ x |(x\in A\wedge x\not\in B)\vee(x\in B\wedge x\not\in A)\}$\\
\{n to m {\em by s}\}& Generate, & 
$\{ x | x=n + is \leq m, i\in\N_0\}$\\
&(default $s = 1$)&\\ 
$\mbox{proj}(A, t)$& Projection $t=(e_1,\ldots,e_n)$& 
The new set will consist of $n$-tuples, with the
$i$-th component beeing the $e_i$-th component of $A$.\\
\hline
\end{tabular}

\medskip
It is possible to write 
* instead of \code{cross}, 
+ instead of \code{union},  
$\backslash$ or - instead of \code{without}
and .. instead of \code{to}.  

\subsubsection{Examples}
{\small 
\begin{verbatim}
set A := { 1, 2, 3 };
set B := { "hi", "ha", "ho" };
set C := { <1,2,"x">, <6,5,"y">, <787,12.6,"oh"> };
set D := A cross B;
set E := { 6 to 9 } union A without { <2>, <3> }; 
set F := <i,j> in Q with i > j and i < 5;
set G := { 1 to 9 } * { 10 to 19 } * { "A", "B" };
set H := proj(G, <3,1>) 
# will give: <"A",1>, <"A",2"> ... <"B",9>
\end{verbatim}
}

\subsection{Parameters}
Parameters can be declared with or without an indexing
set. Without indexing the parameter is just one value, which is either
a number or a string. For indexed parameters there is one
value for each member of the set. There are not much uses for string valued
parameters, so normally they are all numbers. 

Parameters are declared in the following way: 
The keyword \code{param} is followed by the name of the parameter
optionally followed by the indexing set. 
Then after the assignment sign comes a list of pairs. The first element of each
pair is a tuple from the index set, the second element is the value of
the parameter for this index.

\subsubsection{Examples}
{\small 
\begin{verbatim}
param q := 5;
param u[A] := <1> 17, <2> 29, <3> 12;
param w[C] := <1,2,"x"> 1/2, <6,5,"y"> 2/3;
\end{verbatim}
}

\medskip
In the example, no value is given for index $<$787,12.6,"oh"$>$ of
parameter \code{w}, that is assignments need not to be complete. This
is correct as long as it is never referenced.

\subsection{Variables}
Like parameters, variables can be indexed. 
A variable has to be one out of three possible types: 
Continuous (called {\em real}), binary or integer. The default is real. 
Variables may have lower and upper bounds. Defaults are
zero as lower and infinity as upper bound. Binary variables are
always bounded between zero and one.
It is possible to compute the value of the lower or upper bounds
depending on the index for the variable (see last declaration in the example).
\subsubsection{Examples}
{\small 
\begin{verbatim}
var x1;
var x2 binary;
var y[A] real >= 2 <= 18;
var z[<a,b,c> in C] integer 
    >= a * 10 
    <= if b <= 3 then p[b] else 10 end; 
\end{verbatim}
}

\subsection{Objective}
There must be exactly one objective statement in a model. The objective
can be either \code{minimize} or \code{maximize}. Following the
keyword is a name, a colon (\code{:}) and then a term consisting 
of variables.

\subsubsection{Example}
{\small 
\begin{verbatim}
minimize cost: 12 * x1 -4.4 * x2 
   + sum <a> in A : u[a] * y[a]
   + sum <a,b,c> in C with a in E and b > 3 : -a/2 * z[a,b,c];
\end{verbatim}
}

\subsection{Constraints}
The general format for a constraint is 
\code{subto name: term sense term}.
Name can be any name starting with a letter. The term is defined
as in the objective. Sense is one of
\code{<=}, \code{>=} and \code{==}. 
Many constraints can be generated with one statement by the use of the
\code{forall} instruction, see below.

\subsubsection{Examples}
{\small 
\begin{verbatim}
subto time : 3 * x1 + 4 * x2 <= 7;
subto space: sum <a> in A : 2 * u[a] * y[a] >= 50;
subto weird: forall <a> in A : 
   sum <a,b,c> in C : z[a,b,c] == 55;
subto c21: 6 * (sum <i> in A : x[i] 
   + sum <j> in B : y[j]) >= 2;
subto c40: x[1] == a[1] + 
   2 * sum <i> in A do 2*a[i]*x[i]*3 + 4;
\end{verbatim}
}

\subsection{Details on \code{sum} and \code{forall}}
The general forms are

\qquad\code{forall} \emph{index} \code{do} \emph{term}
and 
\code{sum} \emph{index} \code{do} \emph{term}.
 
It is possible to nest several forall instructions.
%It should be noted, that a \code{sum}-expression is a \emph{term}
%itself, so it is possible to nest or concatenate them.
The general form of \emph{index} is 

\qquad\emph{tuple} \code{in} \emph{set} \code{with} \emph{boolean-expression}.

It is allowed to write a colon (\code{:}) instead of \code{do} and a
vertical bar (\code{|}) instead of \code{with}.
The number of components in the \emph{tuple} and in the components of the
mebers of the \emph{set} must match. The \code{with} part of an \emph{index} is
optional. The \emph{set} can be any expression giving a set.

\subsubsection{Examples}
{\small 
\begin{verbatim}
forall <i,j> in X cross { 1 to 5 } without { <2,3> } 
   with i > 5 and j < 2 do 
      sum <i,j,k> in X cross { 1 to 3 } cross Z do 
         p[i] * q[j] * w[j,k] >= if i == 2 then 17 else 53;
\end{verbatim}
}

Note that in the example \emph{i} and \emph{j} are set by the \code{forall}
instruction. So they are fixed for all invocations of \code{sum}.

\subsection{Initialising sets and parameters from a file}
It is possible to load the values for a set or a parameter from a
file. The syntax is

\smallskip
\code{read} \emph{filename} \code{as} \emph{template} 
[\code{skip} \emph{n}] [\code{use} \emph{n}] 
[\code{fs} \emph{s}] [\code{comment} \emph{s}]

\smallskip
\emph{filename} is the name of the file to read.

\emph{template} is a string with a template for the tuples to
generate. Each input line from the file is split in fields. The
splitting is done according to the following rules:
Whenever a space, tab, comma, semicolon or double colon is encountered
a new field is started. Text that is enclosed in double quotes is not
split, the quotes are allways removed. When a field is split all space
and tab charaters around the splitting are removed. If the split is
due to a comma, semicolon or double colon, each occurence of these
characters starts a new field. 

\subsubsection{Examples}
{\small 
All these lines have three fields:
\begin{verbatim}
Hallo;12;3    
Moin   7  2
"Hallo, Peter"; "Nice to meet you" 77
,,2
\end{verbatim}
}

For each component of the tuple the number of the field
to use for the value is given, followed by either a \code{n} if the
field should be interpreted as a number or \code{s} for a string. 
Have a look at the example, it is quite obvious how it works.
After the template some optional modifiers can be given. The order
does not matter. 

\code{skip} \emph{n} instructs to skip the first
\emph{n} lines of the file. 

\code{use} \emph{n} limits the number of
lines to use to \emph{n}. 

\code{comment} \emph{s} sets a list of characters that start
comments in the file. Each line is ended when any of the comment
characters is found.

When a file is read, empty lines are skiped and not counted for the
\code{use} clause. They are counted for the \code{skip} clause.
 
\subsubsection{Examples}
{\small 
\begin{verbatim}
set P := { read "nodes.txt" as <1s> };

nodes.txt:
   Hamburg   ->   <"Hamburg">
   M�nchen   ->   <"M�nchen">
   Berlin    ->   <"Berlin">

set Q := { read "blabla.txt" as "<1s,5n,2n>" skip 1 use 2 };

blabla.txt:
   Name;Nr;X;Y;No     ->   skip   
   Hamburg;12;x;y;7   ->   <"Hamburg",7,12>
   Bremen;4;x;y;5     ->   <"Bremen,5,4>
   Berlin;2;x;y;8     ->   skip
   
param cost[P] := read "cost.txt" as "<1s> 4n" comment "#";

cost.txt:
   # Name Price   ->   skip
   Hamburg 1000   ->   <"Hamburg"> 1000
   M�nchen 1200   ->   <"M�nchen"> 1200
   Berlin  1400   ->   <"Berlin">  1400

param cost[Q] := read "haha.txt" as "<3s,1n,2n> 4s";

haha.txt:
   1:2:ab:con1    ->   <"ab",1,2> "con1"
   2:3:bc:con2    ->   <"bc",2,3> "con1"
   4:5:de:con3    ->   <"de",4,5> "con1"
\end{verbatim}
}  
\newpage

\section{Error messages}
Here are some of the incomprehensible error messages \zimpl
can produce:

\begin{description}
\item[Comparison of different types]\ \\
 It is not possible to compare a number with a string.
\item[xxx of incompatible sets]\ \\
 The members of the two sets involved in operation \emph{xxx} have not
 the same number of components.
\item[Illegal element xxx for symbol]\ \\
 In the initialisation of a parameter, tuple \emph{xxx} is not a member
 in the index set.
\item[WITH not allowed here]\ \\
 When initialising the bounds of variables by index, 
 the \code{with} clause to exclude some of the elements is not allowed.
\item[Comparison of different dimension tuples]\ \\
 Two tuples were compared which have a different number of components.
 This is never a good idea, since such tuples are always different.
\item[Comparison of elements with different types]\ \\
 Two tuples were compared which have a different type for the 
 n-th element. 
 This is never a good idea, since such tuples are always different.
\item[Type error, expected xxx got yyy]\ \\
 In some context type \emph{xxx} was expected, but the interpreter
 found type \emph{yyy}. The meaning of the numbers could be looked up
 in \code{code.h}.
\item[Duplicate constraint name "xxx"]\ \\
 The name given to a \code{subto} statement was allready used with
 another \code{subto} statement. Either change the name, or use 
 \code{-n cm} or \code{-n cf}.
\end{description}

\section{Remarks}
\zimpl is licensed under the GNU general public licence version 2.
For more information on free software see \url{http://www.gnu.org}.
The latest version of \zimpl can be found at
\url{http://www.zib.de/koch/zimpl}.
If you find any bugs you can email me at \url{mailto:koch@zib.de}. Please
include an example that shows the problem.
If somebody extends \zimpl, I am interested in getting patches
to put them back in the main distribution.

\newpage 
\nocite{*}
\bibliographystyle{alpha}
\bibliography{zimpl}

{\vfill\tiny\hfill\fbox{\rcsInfoRevision}}
\end{document}





